\begin{theorem}
\label{theorem-uncountable-nullstellensatz}
Let $k$ be a field. Let $S$ be a $k$-algebra generated over $k$
by the elements $\{x_i\}_{i \in I}$. Assume the cardinality of $I$
is smaller than the cardinality of $k$. Then
\begin{enumerate}
\item for all maximal ideals $\mathfrak m \subset S$ the field
extension $\kappa(\mathfrak m)/k$
is algebraic, and
\item $S$ is a Jacobson ring.
\end{enumerate}
\end{theorem}

\begin{proof}
If $I$ is finite then the result follows from the Hilbert Nullstellensatz,
Theorem \ref{theorem-nullstellensatz}. In the rest of the proof we assume
$I$ is infinite. It suffices to prove the result for
$\mathfrak m \subset k[\{x_i\}_{i \in I}]$ maximal in the polynomial
ring on variables $x_i$, since $S$ is a quotient of this.
As $I$ is infinite the set
of monomials $x_{i_1}^{e_1} \ldots x_{i_r}^{e_r}$, $i_1, \ldots, i_r \in I$
and $e_1, \ldots, e_r \geq 0$ has cardinality at most equal to the
cardinality of $I$. Because the cardinality of $I \times \ldots \times I$
is the cardinality of $I$, and also the cardinality of
$\bigcup_{n \geq 0} I^n$ has the same cardinality.
(If $I$ is finite, then this is not true and
in that case this proof only works if $k$ is uncountable.)

\medskip\noindent
To arrive at a contradiction pick $T \in \kappa(\mathfrak m)$ transcendental
over $k$. Note that the $k$-linear map
$T : \kappa(\mathfrak m) \to \kappa(\mathfrak m)$
given by multiplication by $T$ has the property that $P(T)$ is invertible
for all monic polynomials $P(t) \in k[t]$.
Also, $\kappa(\mathfrak m)$ has dimension at most the cardinality of $I$
over $k$ since it is a quotient of the vector space
$k[\{x_i\}_{i \in I}]$ over $k$ (whose dimension is $\# I$ as we saw above).
This is impossible by Lemma \ref{lemma-dimension}.

\medskip\noindent
To show that $S$ is Jacobson we argue as follows. If not then
there exists a prime $\mathfrak q \subset S$ and an element $f \in S$,
$f \not \in \mathfrak q$ such that $\mathfrak q$ is not maximal
and $(S/\mathfrak q)_f$ is a field, see
Lemma \ref{lemma-characterize-jacobson}.
But note that $(S/\mathfrak q)_f$ is generated by at most $\# I + 1$ elements.
Hence the field extension $(S/\mathfrak q)_f/k$ is algebraic
(by the first part of the proof).
This implies that $\kappa(\mathfrak q)$ is an algebraic extension of $k$
hence $\mathfrak q$ is maximal by
Lemma \ref{lemma-finite-residue-extension-closed}. This contradiction
finishes the proof.
\end{proof}